\chapter{Intoducción}

\section{Motivación y Planteamiento}
Según los datos de la Organización Mundial de la Salud, más de $1.500$ millones de personas en el mundo viven con algún grado de pérdida auditiva, de las cuales aproximadamente $430$ millones requieren atención especializada. Para una parte significativa de esta población, la lengua de signos es su principal medio de comunicación natural. Sin embargo, el acceso a servicios esenciales como lo son la sanidad, la educación o la administración pública sigue dependiendo casi exclusivamente de intérpretes humanos, un recurso escaso, costoso y con disponibilidad geográfica muy desigual. Esta dependencia limita la autonomía de las personas sordas y perpetúa una brecha de inclusión social difícil de cerrar sin el apoyo de soluciones tecnológicas escalables.

En los últimos años ha habido enormes avances en las técnicas de \textit{Procesamiento del Lenguaje Natural}, \textit{Visión por Computador} y \textit{Aprendizaje Profundo}, así como en la eficiencia y capacidad computacional de los equipos, sin embargo las lenguas de signos han recibido una atención comparativamente menor, y los avances no han alcanzado el nivel de madurez ni la cobertura lingüística que sí presentan los sistemas para lenguas orales. Los sistemas de traducción automática, los asistentes conversacionales y las interfaces de voz han transformado la forma en que las personas oyentes interactúan con la tecnología, pero siguen siendo inaccesibles para quienes se comunican mediante el canal viso-gestual.

Para abordar esta brecha, la investigación en traducción automática de lengua de signos se articula principalmente en torno a dos direcciones complementarias: por un lado, la traducción de texto o voz a lengua de signos, \textbf{Text2Sign}, orientada a facilitar la comprensión por parte de personas sordas; por otro, la traducción de lengua de signos a texto, \textbf{Sign2Text}, cuyo objetivo es permitir que las personas oyentes o los sistemas automáticos comprendan los mensajes signados. El presente trabajo se centra en esta segunda dirección, \textbf{Sign2Text}, por su impacto directo en la autonomía comunicativa de las personas sordas y por los retos técnicos abiertos que aún presenta.

La tarea \textbf{Sign2Text}, también conocida como \textit{Sign Language Recognition and Translation}, persigue la generación automática de texto a partir de una secuencia de signos capturada en vídeo. Un sistema de este tipo permitiría subtitular en tiempo real una conversación signada, dotar a los asistentes virtuales de la capacidad de entender a usuarios sordos, o facilitar el acceso a servicios públicos sin necesidad de un intérprete presente. En definitiva, actuaría como un puente de comunicación directo entre la comunidad sorda y los sistemas digitales diseñados, hasta ahora, pensando exclusivamente en usuarios oyentes.

Desde el punto de vista técnico, Sign2Text es un problema especialmente complejo. A diferencia del reconocimiento de voz, donde la señal de entrada es unidimensional y continua, la lengua de signos es inherentemente multimodal: la información se distribuye simultáneamente entre la configuración y el movimiento de las manos, la expresión facial, la postura corporal y el uso del espacio. Esta riqueza expresiva, que dota a las lenguas de signos de toda su potencia comunicativa, supone al mismo tiempo uno de los mayores desafíos para su modelado automático. A esto se le suma la gran variedad de lenguas de signos que existen, cada una con sus gestos y expresiones visuales únicas, además en la mayoria de los casos, la lengua de signos de una determinada región carece de relación gramatical con la lengua natural de dicha región. Esto dificulta significativamente la comprensión y traducción directa entre un lenguaje natural y de signos.

Además de esta complejidad multimodal, la naturaleza visual de la tarea presenta consideraciones relevantes en materia de privacidad. Procesar vídeos de personas requiere garantías sobre el tratamiento de datos personales, especialmente en entornos sensibles como la sanidad o la educación. Por ello, un enfoque que minimice la exposición de datos desde el origen resulta deseable. A raiz de esto nace la traducción de signos basados en \textit{keypoints}, esto es puntos predeterminados en el cuerpo humano que resumen la información de la postura, expresión e información de manos, capturando la información gestual relevante sin preservar la identidad visual de la persona, permitiendo así un procesamiento más ligero, portable y respetuoso con la privacidad. De forma paralela, el uso de este enfoque, reduce drásticamente la dimensionalidad de los datos, optimizando los costes computacionales de entrenamiento e inferencia y facilitando su despliegue en equipos menos potentes.

\section{Objetivo}
